\SourceLeaf{365}%
cases it leads to the goal more directly and by a shorter route than analysis, especially when analysis is applied without dexterity. It then has an elegance of its own. Beyond this, it will remain indispensable, especially for formal training and for study at an early age. It guards against a one-sided approach, sharpens the sense of rigor and clarity, and gives understanding a vividness and immediacy that analytical methods do much less to foster and sometimes even put at risk. For these reasons it is a pleasure to see that in recent decades some French geometers have begun to give particular attention to the branch of geometry that deals with relations among points and lines that do not lie in a single plane, with the relations of different planes to one another, and with lines of double curvature and curved surfaces. To the extent that they use geometrical methods alone, they cultivate it as a separate discipline under the name \textit{Géométrie descriptive}. We must give this book on the subject particular credit for the clarity and economy of its exposition, the careful progression from easier to harder matters, and the wealth of fresh perspectives and successful developments. We therefore recommend studying it as sustaining food for the mind, which can surely do much to revive and preserve the genuine spirit of geometry that is sometimes missed elsewhere in modern mathematics. Beyond the purely scientific side of these inquiries, there are the many uses they have in arts concerned with spatial relations, especially drawing, land surveying, architecture and fortification. In this respect, too, the author has managed to make his work more interesting through various applications, though for the most part he has suggested them rather than carried them out.
